This dissertation was prepared with the assistance of generative AI tools, and this statement describes how they were used.

\paragraph{Writing assistance.}
The chapters of this thesis are based on the author's first-author publications listed in \Cref{intro:sec:publications}.
An AI coding assistant (Claude Code, Anthropic) was used to convert the \LaTeX{} sources of those papers into thesis chapters, to draft the background, introduction, and conclusion chapters from the content of the papers and the author's thesis proposal, and to assist with editing, formatting, and bibliography management.
All text produced with AI assistance was reviewed and revised by the author, who takes full responsibility for the content, claims, and any errors in this document.
The research contributions, experimental results, and conclusions are those of the author and collaborators as reported in the underlying publications.

\paragraph{AI systems as research artifacts.}
\Cref{ch:fastkernels,ch:kernelseek} study large language model agents as their subject.
The GPU kernels evaluated in those chapters, and in \Cref{ch:kernelseek} the kernels themselves, the baseline agent, the evaluation harness, and the trace-parsing tools, were produced by LLM-based coding agents as part of the research, as each chapter states.
The design of the experiments, the analysis of the results, and the claims drawn from them are the author's.

\paragraph{Scope.}
No generative AI tool was used to produce or alter the experimental data reported in this thesis beyond what the underlying publications disclose.
